\chapter{Child Abuse and Neglect}

\section*{Definition and Overview}
Child abuse, also known as child maltreatment, is any act by a parent, carer, or other adult in a position of trust that results in harm to a child, or that creates a substantial risk of harm. It takes several forms: physical abuse, emotional (psychological) abuse, sexual abuse, and neglect, which is the failure to meet a child's basic needs for food, shelter, supervision, medical care, or affection. Children can also be harmed by witnessing family violence. Abuse affects children of every age, background, and community. It is a serious public health issue with lifelong consequences, and healthcare professionals have a duty to recognise its signs and to act to protect the child. This chapter explains the types of abuse, warning signs, and the roles of assessment and protection.

\section*{Epidemiology}
Child maltreatment is a global problem. It is estimated that a substantial proportion of children worldwide experience physical, emotional, or sexual abuse or neglect, although the true figures are uncertain because much abuse goes unreported. Risk is highest in the very young, in children with disabilities, and in families affected by poverty, domestic violence, parental substance misuse, or mental illness. Boys and girls are abused in similar overall numbers, although patterns differ by type: girls are more often victims of sexual abuse, while boys experience more physical abuse. Most abuse is perpetrated by parents or close carers, not by strangers. Reporting rates have risen in many countries as awareness and child protection systems have improved.

\section*{Aetiology and Pathophysiology}
\medskip
\textbf{Forms of maltreatment:}
\begin{itemize}
    \item \textbf{Physical abuse:} hitting, shaking, throwing, burning, or otherwise injuring a child deliberately
    \item \textbf{Emotional abuse:} persistent humiliation, threats, rejection, or exposure to family violence that damages a child's development and self-worth
    \item \textbf{Sexual abuse:} involving a child in sexual activity, including touching, exploitation, or exposing them to sexual material
    \item \textbf{Neglect:} the ongoing failure to provide adequate food, clothing, shelter, supervision, medical care, or affection
    \item \textbf{Medical child abuse:} a carer fabricating or inducing illness in a child, causing unnecessary investigations and treatment
\end{itemize}

\medskip
\textbf{Risk factors.} Abuse results from a combination of child, parent, and environmental factors:
\begin{itemize}
    \item \textbf{Child factors:} young age, disability, and difficult temperament
    \item \textbf{Parent factors:} substance misuse, mental illness, a history of being abused, and poor parenting skills
    \item \textbf{Environmental factors:} poverty, social isolation, domestic violence, and community violence
    \item \textbf{Triggers:} stressful events and, particularly in infancy, the distress caused by prolonged crying
\end{itemize}

\section*{Clinical Features}
\begin{itemize}
    \item \textbf{Bruises:} in unusual places such as the face, ears, neck, buttocks, or abdomen, in clusters, or in patterns matching objects; bruising in babies who cannot yet crawl or cruise is always concerning
    \item \textbf{Other injuries:} burns in patterns (for example, scald lines or cigarette burns), bite marks, and fractures in a child too young to fall or move independently
    \item \textbf{Behavioural signs:} sudden changes in behaviour, fear of a particular adult, flinching, withdrawal, or being unusually watchful
    \item \textbf{Neglect:} poor growth, dirty or inadequate clothing, untreated medical or dental problems, and frequent hunger
    \item \textbf{Emotional signs:} very low self-esteem, anxiety, depression, and self-harm in older children
    \item \textbf{Sexual abuse signs:} genital injury, sexually transmitted infection, inappropriate sexual knowledge or behaviour for the child's age
    \item \textbf{Disclosure:} a child describing abuse is the strongest signal and must always be taken seriously
\end{itemize}

\section*{Differential Diagnosis}
\begin{itemize}
    \item Accidental injury, including falls and everyday bruises on bony prominences such as shins and elbows
    \item Bleeding disorders such as haemophilia, or conditions such as immune thrombocytopenia that cause easy bruising
    \item Skin conditions including Mongolian spots, birthmarks, and eczema
    \item Bone conditions such as osteogenesis imperfecta, which causes fractures with minor trauma
    \item Cultural practices such as cupping or coining that leave marks
    \item Behavioural and mental health conditions mistaken for emotional abuse
\end{itemize}

\section*{When to Seek Medical Care}
Anyone who suspects a child is being abused or neglected should report it. Healthcare professionals have clear duties to notify child protection services. Family members, neighbours, and teachers can also report concerns confidentially.

\medskip
\textbf{Call 000 (or local emergency services) for:}
\begin{itemize}
    \item A child with a serious or unexplained injury, or who appears seriously unwell
    \item A child who is unconscious, drowsy, or vomiting after an injury
    \item A child in immediate danger from a violent or threatening situation
    \item Suspected sexual assault that has just occurred
\end{itemize}

\section*{Diagnosis and Investigations}
\begin{itemize}
    \item \textbf{History:} a careful account of the injury or concern, noting whether the explanation matches the child's age, development, and the nature of the injury
    \item \textbf{Examination:} a full, gentle examination of the child, carefully documenting all injuries, bruises, and marks
    \item \textbf{Growth and developmental assessment:} weight, height, and head circumference, with attention to failure to thrive
    \item \textbf{Imaging:} a skeletal survey (X-rays of the whole skeleton) is performed in children under two years when physical abuse is suspected, to look for old or hidden fractures
    \item \textbf{Brain imaging:} CT or MRI may be needed if abusive head trauma or shaken baby syndrome is suspected
    \item \textbf{Blood tests:} clotting studies and other tests to exclude medical conditions that mimic abuse
    \item \textbf{Forensic assessment:} specialised medical examination and documentation in cases of suspected sexual abuse
    \item \textbf{Child protection referral:} notification to child protection services and, where appropriate, police involvement, is part of the standard response
\end{itemize}

\section*{Management}
\begin{itemize}
    \item \textbf{Child safety first:} the immediate priority is to ensure the child is safe, which may mean admission to hospital while concerns are investigated
    \item \textbf{Multidisciplinary assessment:} a coordinated assessment by paediatricians, child protection workers, and police where needed
    \item \textbf{Treatment of injuries:} medical and surgical treatment of any physical injuries
    \item \textbf{Emotional support:} psychological support and counselling for the child and appropriate family members
    \item \textbf{Family support services:} parenting programs, substance misuse treatment, and mental health care for parents when the family can be safely supported
    \item \textbf{Child protection plan:} ongoing case management by statutory services, with regular reviews of the child's safety
    \item \textbf{Legal processes:} involvement of the courts when care arrangements need to be made or changed
\end{itemize}

\section*{Complications}
\begin{itemize}
    \item Serious physical injury, including brain injury, fractures, and burns
    \item Permanent disability and, in severe cases, death
    \item Post-traumatic stress, anxiety, and depression
    \item Long-term mental health problems, including self-harm and substance misuse in adolescence and adulthood
    \item Developmental delay and poor educational outcomes
    \item Difficulties forming trusting relationships and increased risk of later victimisation or perpetration
    \item Physical health consequences in adult life, including cardiovascular and metabolic disease linked to early stress
\end{itemize}

\section*{Prognosis and Outcomes}
With early recognition and effective intervention, many children who have experienced abuse recover well, particularly when they receive stable, nurturing care and appropriate treatment. Children who remain in abusive environments, or who have experienced severe or repeated maltreatment, are at higher risk of lasting psychological and physical harm. The outlook depends heavily on the swiftness of protection, the quality of subsequent care, and the support available to the child and family. Prevention and early help are the most effective ways to improve outcomes.

\section*{Prevention and Self-Care}
\begin{itemize}
    \item Support parents with positive parenting programs and information about child development
    \item Encourage calm responses to infant crying and safe strategies such as placing the baby down and taking a break
    \item Provide respite and social support for stressed or isolated families
    \item Address parental substance misuse, mental illness, and domestic violence through services
    \item Teach children age-appropriate safety skills about their bodies and consent
    \item Ensure that professionals and community members know how to report concerns
    \item Strengthen community programs such as home visiting for families at risk
\end{itemize}

\section*{Sources and Further Reading}
The following reputable sources provide further information on child abuse and neglect:
\begin{itemize}
    \item World Health Organization. \textit{Child maltreatment fact sheets and prevention resources.} https://www.who.int/
    \item Centers for Disease Control and Prevention. \textit{Child abuse and neglect prevention resources.} https://www.cdc.gov/
    \item NHS. \textit{Health A--Z and safeguarding children information.} https://www.nhs.uk/conditions/
    \item NICE. \textit{Guidelines on child abuse and neglect.} https://www.nice.org.uk/
    \item MedlinePlus. \textit{Health topics on child safety and wellness.} https://medlineplus.gov/
    \item NSPCC. \textit{Safeguarding children and reporting concerns.} https://www.nspcc.org.uk/
\end{itemize}
